\chapter{Experimental Setup}
\label{chap:experimental-setup}

Once the task formulations are fixed, the main experimental risk is no longer
architectural ambiguity but unfair comparison. Models can look better or worse
simply because they were trained on different splits, selected with different
criteria, or evaluated through different switch rules. The protocol below makes
these choices explicit. All experiments use the same raw signal sources and the
same external-holdout principle, while each task still produces its own
supervised dataset because the target variable and the set of valid decision
timestamps are different.

The setup follows three constraints. First, the test dataset is held out as a
complete external file and is never used for model selection. Second, all fitted
transformations, such as feature standardization, are estimated on the training
split only. Third, all final operational comparisons are made after converting
model outputs into binary switch decisions and applying the shared
post-processing rule defined in \cref{sec:shared-switch-post-processing}.

\section{Dataset}
\label{sec:experimental-dataset-configuration}

The common operational threshold used in the final experiments is
\(\theta=\SI{10.0}{dB}\). Unless stated otherwise, samples are represented on a
\SI{30}{s} grid and the decision context contains the last \(L=30\) signal
observations, corresponding to approximately \SI{15}{min} of history. The
operational persistence horizon used for switch decisions is
\(T_{\min}=\SI{300}{s}\), i.e. \SI{5}{min}. This horizon is used as the forecast
horizon interpretation for autoregressive switch conversion, as the duration
threshold for current-level persistence, as the minimum long-fade duration for
long-fade detection, and as the survival horizon for survival persistence.

The external test source is the complete file
\texttt{fc-uplink-fade.csv}. It is excluded from all training, validation, grid
search, early stopping, and model-selection steps. The remaining signal files
form the development pool. Within the development pool, train and validation
splits are constructed chronologically at event level when event identities are
available, so that strongly related windows from the same event are not split
between training and validation.

\subsection{Autoregressive Forecasting Dataset}

The autoregressive dataset uses past windows and future signal trajectories.
The main configuration uses \(L=30\) and prediction length \(h=10\), so each
target covers the next \SI{300}{s}. The resulting external-holdout dataset
contains 19,412 training windows, 4,783 validation windows, and 7,708 test
windows under the inclusive \(S_i\ge\SI{10.0}{dB}\) threshold convention. A
separate diagnostic Chronos experiment also uses \(L=120\) while keeping
\(h=10\); this longer context is treated as an additional Chronos run, not as
the common autoregressive configuration for all models.

\subsection{Current-Level Persistence Dataset}

The current-level persistence dataset uses the same \(L=30\) context length and
\(\delta=\SI{0.5}{dB}\). Its target is the residual time for which the future
signal remains above \(S_i-\delta\). Since this target requires observing a
future recovery below a sample-specific level, a small number of candidate
windows is right-censored and excluded from ordinary duration-regression
training. The final uncensored regression dataset contains 19,598 training
windows, 4,798 validation windows, and 7,692 test windows. The excluded
censored candidate windows account for approximately 2.12\%, 2.74\%, and
3.22\% of the train, validation, and test candidate windows respectively.

\subsection{Long-Fade Detection Dataset}

The long-fade detection dataset is built only inside grouped fade events. It
uses \(L=30\), \(\theta=\SI{10.0}{dB}\), and a long-fade label defined by event
duration at least \(T_{\min}=\SI{300}{s}\). The resulting dataset contains 4,589
training samples from 62 grouped events, 790 validation samples from 15 grouped
events, and 2,569 test samples from 25 grouped events. This dataset is much
smaller than the autoregressive and current-level datasets because it does not
generate samples outside grouped fade events.

\subsection{Survival-Persistence Dataset}

The survival-persistence dataset also uses \(L=30\), but keeps censored samples
instead of discarding them. Each sample is generated inside a detected fade
event and stores lower and upper time bounds for the residual event duration.
Observed samples have equal lower and upper bounds, while censored samples have
an infinite upper bound. The final survival dataset contains 1,729 training
samples, 729 validation samples, and 921 test samples. The validation split
contains 6 censored samples from one event, while the external test split
contains 70 censored samples from two events.

\subsection{Small-Gap Imputation Policy}

Small-gap imputation is enabled only through the shared conservative policy.
It may fill short gaps up to \SI{90}{s}, corresponding to at most two inserted
\SI{30}{s} samples, by linear interpolation in time. The policy is applied
before task-specific dataset construction where explicitly configured, but it
does not bridge acquisition-segment boundaries or long missing intervals.

\section{Train, Validation, and Test Protocol}
\label{sec:train-validation-test-protocol}

The experiments use a three-way protocol: training, validation, and external
test. The training split is used to fit model parameters. The validation split is
used for early stopping, hyperparameter selection, architecture selection, and
any comparison among alternative variants of the same task. The external test
split is used only once a model or configuration has already been selected.

The distinction between validation and test is especially important in this
work, because the models are evaluated not only through their native predictive
losses but also through derived switch decisions. A threshold or architecture
that looks attractive on the test set would no longer provide an unbiased
estimate of future operational behavior. For this reason, the final selected
models are chosen using validation-only criteria. Exploratory operating points
reported outside this rule are explicitly treated as diagnostics rather than as
calibrated deployment thresholds.

For trainable neural models, early stopping or fixed-epoch training is governed
by validation-only information during the selection phase. After the best
configuration is chosen, the published run is either retrained on the full
development set, i.e. training plus validation, or the selected validation-stage
checkpoint is promoted. For reproducibility, this finalization choice is not
left implicit: each published run stores in its metadata whether the test
artifact was obtained from a train-validation refit or from a promoted
validation-stage checkpoint, together with the selected epoch, configuration,
random seed, and dataset fingerprint when available. The external test split is
evaluated only after the selection step has been completed. XGBoost models
follow the same principle: validation selects the hyperparameters, and the
selected configuration is then published before test evaluation. Chronos is the
only exception, because it is evaluated in zero-shot mode and therefore loads no
training or validation data.

All splits preserve the signal-only policy. Time, event identifiers, dataset
identifiers, and quality flags are stored as metadata for alignment,
diagnostics, and grouping, but they are not used as model inputs unless a
feature is explicitly described as an online scalar derived from the past signal
window or from the current event state available at the decision time.

\section{Hyperparameter Search}
\label{sec:hyperparameter-search}

Hyperparameter search is performed separately for each task and model family.
Each search is defined by a YAML configuration file and produces a structured
trial table, selected-parameter file, final run folder, model artifact folder,
and metadata file. The grid-search runners assign independent trials across all
available CUDA devices when possible. If CUDA is unavailable, the code falls
back to MPS when supported, and finally to CPU.

\begin{table}[htbp]
  \centering
  \caption{Summary of the main model-selection searches.}
  \label{tab:model-selection-searches}
  \small
  \setlength{\tabcolsep}{4pt}
  \begin{tabular}{
    @{}
    >{\raggedright\arraybackslash}p{0.18\textwidth}
    >{\raggedright\arraybackslash}p{0.36\textwidth}
    >{\raggedright\arraybackslash}p{0.11\textwidth}
    >{\raggedright\arraybackslash}p{0.25\textwidth}
    @{}
  }
    \toprule
    Task & Model family & Trials & Selection metric \\
    \midrule
    Autoregressive & GRU S2V, raw/context-standard & 72 each & validation RMSE \\
    Autoregressive & GRU encoder-decoder, raw/context-standard & 216 each & validation RMSE \\
    Autoregressive & PatchTST, raw/context-standard & 864 each & validation RMSE \\
    Autoregressive & XGBoost, raw/context-standard & 1152 each & validation RMSE \\
    Current-level persistence & Shapelet MLP/Transformer/Convolution & 216 each & validation MAE \\
    Current-level persistence & XGBoost duration regressor & 1152 & validation MAE \\
    Long-fade detection & XGBoost/TCN/Shapelet-Convolution classifiers & 8928 total & validation AUPRC \\
    Survival persistence & XGBoost-AFT & 1920 & validation AFT negative log-likelihood \\
    Survival persistence & Discrete-time TCN & 1620 & validation Integrated Brier Score \\
    \bottomrule
  \end{tabular}
\end{table}

The search dimensions and validation-only selection rules are summarized in
\cref{tab:model-selection-dimensions}. All error-based selection metrics are
computed after returning predictions to their reported physical scale.

\begin{table}[htbp]
  \centering
  \caption{Main hyperparameter dimensions and validation-only selection rules.}
  \label{tab:model-selection-dimensions}
  \small
  \setlength{\tabcolsep}{4pt}
  \renewcommand{\arraystretch}{1.25}
  \begin{tabular}{
    @{}
    >{\raggedright\arraybackslash}p{0.24\textwidth}
    >{\raggedright\arraybackslash}p{0.47\textwidth}
    >{\raggedright\arraybackslash}p{0.20\textwidth}
    @{}
  }
    \toprule
    Task and model family & Tuned dimensions & Selection rule \\
    \midrule
    Autoregressive GRU
      & input variant, recurrent capacity, number of layers, dropout, learning rate, weight decay, teacher forcing for encoder-decoder
      & validation RMSE in raw signal scale \\
    Autoregressive PatchTST
      & input variant, patch length and stride, Transformer width, heads, layers, dropout, learning rate, weight decay
      & validation RMSE in raw signal scale \\
    Autoregressive XGBoost
      & input variant, tree depth, estimators, learning rate, subsampling, column sampling, child weight, regularization
      & validation RMSE in raw signal scale \\
    Current-level persistence
      & shapelet representation size, neural head capacity, scalar-context embedding, dropout, learning rate; XGBoost tree parameters
      & validation MAE in seconds \\
    Long-fade detection
      & classifier family, sequence or lag representation, model capacity, tree or neural regularization
      & validation AUPRC \\
    Survival persistence
      & AFT tree parameters; TCN binning, input representation, sample weighting, architecture scale, seed
      & AFT negative log-likelihood or Integrated Brier Score \\
    \bottomrule
  \end{tabular}
\end{table}

\section{Forecast, Regression, Classification, and Survival Metrics}
\label{sec:native-task-metrics}

The metric definitions are introduced in
\cref{sec:background-operational-evaluation}. This section specifies how those
metrics are used in the experimental protocol.

For autoregressive forecasting, MAE and RMSE are computed after all predictions
have been returned to raw signal scale. They are reported globally and by
horizon step. Validation RMSE is used for neural and XGBoost model selection,
because large forecast errors near the operational threshold can strongly affect
the derived switch decision.

For current-level persistence, the models are trained on
\(\log(1+D_i^{\text{CLP}})\), but native evaluation is performed after mapping
predictions back to seconds. Validation MAE in seconds is used for model
selection, and test MAE/RMSE in seconds are retained as diagnostic duration
metrics.

For long-fade detection, the native output is a probability that the grouped
event is a long fade. Validation AUPRC is used for model selection because the
classification problem is imbalanced. AUROC and thresholded classification
metrics are reported as secondary diagnostics.

For survival persistence, XGBoost-AFT is selected by validation AFT negative
log-likelihood, while the discrete-time TCN is selected by validation
Integrated Brier Score. Concordance index, Brier scores, and survival
likelihood values are reported to diagnose whether the predicted residual
duration distribution is meaningful before it is converted into a switch.

\section{Switch Metrics}
\label{sec:switch-metrics}

The switch metrics are also defined in
\cref{sec:background-operational-evaluation}. In the experiments, they are
computed after task-specific raw decisions have been converted into the common
post-processed switch sequence. The evaluated sequence is always compared with
the task-scoped Perfect Switch on the same timestamp grid.

The reported switch tables include pointwise metrics, active switch duration,
the number of switch islands, detected and missed Perfect Switch events, false
switch events, onset delay when defined, and the amount of change introduced by
post-processing. These quantities are reported because the operational cost of
backup-link activation depends on the duration and fragmentation of switch
periods, not only on pointwise classification agreement.

\section{Use of Shared Switch Post-Processing}
\label{sec:experimental-switch-post-processing}

The post-processing rule itself is defined with the reference framework in
\cref{sec:shared-switch-post-processing}. In the experiments, the same
configured rule is applied unchanged after each task-specific raw switch
conversion. The primary evaluated sequence is \texttt{model\_switch}, which is
the held minimum-time sequence introduced in
\cref{sec:shared-switch-post-processing}. The auxiliary diagnostic columns are:

\begin{itemize}
  \item \texttt{model\_switch\_raw}, the immediate task-specific decision;
  \item \texttt{model\_switch\_min\_island\_old}, after minimum-duration
        extension;
  \item \texttt{model\_switch\_min\_time}, after the above-threshold holding
        rule;
  \item \texttt{model\_switch\_adjusted}, the legacy adjusted diagnostic
        variant.
\end{itemize}

Only \texttt{model\_switch} is used as the main operational sequence; the other
columns are retained to audit how much post-processing changes the immediate
model decision.

\section{Intra-Task Comparison Protocol}
\label{sec:intra-task-comparison-protocol}

Intra-task comparisons evaluate models that solve the same supervised problem.
They therefore share the same dataset, split, native target, Perfect Switch
grid, and switch-conversion rule. This is the cleanest comparison level because
the only intended difference is the model family or input representation.

For autoregressive forecasting, the comparison includes GRU
sequence-to-vector, GRU encoder-decoder, PatchTST, XGBoost, and Chronos
zero-shot runs. Forecast metrics are compared where available, and switch
metrics are computed after applying the autoregressive conversion rule defined
in \cref{chap:task-formulations}.

For current-level persistence, the comparison includes the three learnable
shapelet regressors and the XGBoost duration regressor. The models are selected
by validation duration error and then evaluated through the common
duration-to-switch conversion.

For long-fade detection, the comparison includes XGBoost, TCN, and
shapelet-convolution classifiers. The models are selected by validation AUPRC
and then evaluated through the common probability-to-switch conversion.

For survival persistence, XGBoost-AFT and the grid-selected discrete-time TCN
are compared through survival metrics and through the survival-to-switch
conversion at the configured operational horizon.

Each intra-task comparison also produces event-level plots. The plots show the
raw signal, the operational threshold, the Perfect Switch, and the
post-processed model switch sequence. For long-fade and survival models, the
visual window is extended to 90 minutes before and after the event timestamp,
while switches outside the task-native decision grid are displayed as zero. This
makes event plots easier to compare visually without changing the metric
denominator.

\section{Cross-Task Comparison Protocol}
\label{sec:cross-task-comparison-protocol}

Cross-task comparison is more difficult than intra-task comparison because the
tasks do not necessarily produce decisions at the same timestamps. The
autoregressive task produces one decision for each valid forecast window.
Current-level persistence produces decisions for valid duration-regression
samples. Long-fade detection and survival persistence produce decisions only
inside their own event definitions. Directly comparing only the common
intersection of all timestamps would hide this difference in coverage.

The primary cross-task protocol therefore aligns all methods to a common
reference grid derived from the external-holdout Perfect Switch. If a method has
no native decision at a reference-grid timestamp, the missing value is filled as
zero. This denominator is called
\texttt{reference\_grid\_missing\_as\_zero}. It evaluates both the correctness of
the decisions that a method makes and the operational consequence of not
producing decisions at some reference-grid timestamps.

A second denominator,
\texttt{strict\_common\_timestamp\_intersection}, is also computed as a
diagnostic. It keeps only timestamps where all compared methods have native
decisions. This view is useful for asking which method behaves better when all
methods are active on the same timestamps, but it is not sufficient as the main
comparison because it removes the coverage problem.

The cross-task analysis therefore reports both switch metrics and native
decision coverage. Coverage is the fraction of the reference grid on which a
method actually produced a native decision before missing values were filled.
This is essential for interpreting models such as long-fade and survival
methods, whose task definitions intentionally restrict the decision set.

\section{Hardware and Reproducibility}
\label{sec:hardware-reproducibility}

The implementation is based on Python. Neural models are implemented with
PyTorch, while tree-based baselines use XGBoost. Experiment configuration is
stored in YAML files, and generated artifacts are organized by task, selection
identifier, and run identifier.

The source-code repository is linked in \cref{sec:introduction-contributions}.
Its structure and artifact schemas are described in
\cref{app:implementation-details}.

All model runners use centralized device selection. The preferred device order
is CUDA, then Apple MPS when supported, then CPU. Grid searches consist of
independent trials and therefore distribute trials across all available CUDA
devices, for example \texttt{cuda:0}, \texttt{cuda:1}, and \texttt{cuda:2}, when
the machine exposes multiple GPUs. Final run metadata records the selected
device, configuration fingerprint, dataset fingerprint when available, model
family, architecture, variant, seed, and output paths.

The main random seed is 42 unless otherwise stated. Neural training stores the
best checkpoint and training metadata. XGBoost runs store the fitted model and
the resolved configuration. Chronos stores run metadata and predictions but no
trained checkpoint, because no project-specific training is performed. All
published runs save prediction files so that switch comparisons can be
regenerated without retraining the models.

The result structure separates four artifact types. First, dataset-preparation
artifacts store split summaries, metadata, and quality checks. Second,
grid-search artifacts store trial tables and selected configurations. Third,
run artifacts store model predictions, native metrics, plots, and run metadata.
Fourth, comparison artifacts store intra-task switch evaluations, model
summaries, and cross-task switch plots. This separation is used to make later
thesis tables reproducible from saved outputs rather than from manually edited
numbers.
